% ---------------------------------------------------------------------
\section{Les acteurs de l’impression à la demande}
% ---------------------------------------------------------------------

On peut classer les acteurs en trois familles : les \textbf{plateformes}
(qui impriment et distribuent pour vous), les \textbf{imprimeurs}
(qui fabriquent, à vous de vendre) et l’\textbf{interprofession}
(les bases de données et institutions incontournables).

\subsection{Les plateformes d’impression et de distribution}

\begin{fiche}{Amazon KDP}{kdp.amazon.com}
\textbf{Le géant incontournable de la vente en ligne.} Mise en ligne
gratuite, ISBN gratuit possible, livres brochés et reliés, ebook Kindle
dans la même interface.\par
\textbf{Points forts :} gratuit, simple, accès direct au premier
libraire en ligne de France, redevances lisibles.\par
\textbf{Limites :} vente quasi exclusivement sur Amazon ; quasiment
invisible pour les librairies françaises ; qualité d’impression
inégale selon les sites ; la « distribution étendue » vise surtout les
marchés anglophones.
\end{fiche}

\begin{fiche}{BoD — Books on Demand}{bod.fr}
\textbf{Pionnier européen de l’IAD}, d’origine allemande (Norderstedt,
près de Hambourg), présent en France depuis de nombreuses années avec une
équipe francophone. BoD s’adresse à la fois aux \textbf{auteurs
indépendants} et aux \textbf{maisons d’édition}.\par
\textbf{Ce que fait BoD :} impression à la demande (principalement en
Allemagne), attribution d’ISBN, référencement et distribution auprès des
libraires en ligne (Amazon, Fnac, Decitre, Cultura…) et des
\textbf{librairies françaises via Dilicom}, en Allemagne et à
l’international ; distribution des ebooks.\par
\textbf{Points forts :} un des rares services qui rendent un livre
autoédité \textbf{commandable en librairie} en France ; contrat non
exclusif, l’auteur garde ses droits ; offres dédiées aux éditeurs
(utilisation de vos propres ISBN et de votre nom d’éditeur) ;
interlocuteurs en français.\par
\textbf{Limites :} selon l’offre choisie, frais de mise en ligne et/ou
frais annuels ; délais de livraison un peu plus longs en France
(impression en Allemagne) ; marge par exemplaire à calculer
soigneusement avec la remise libraire.\par
\textbf{Pour qui ?} L’auteur qui veut être présent en librairie et en
ligne avec un seul interlocuteur ; la petite maison qui veut maintenir
son fonds disponible sans stock.
\end{fiche}

\begin{fiche}{IngramSpark}{ingramspark.com}
\textbf{Le réseau mondial d’Ingram} (premier distributeur de livres aux
États-Unis), via son imprimeur Lightning Source, qui dispose d’un site
d’impression en France, à Maurepas (Yvelines), exploité en partenariat
avec Hachette Livre.\par
\textbf{Points forts :} plus de frais de mise en ligne ; large réseau
de distribution international ; remise libraire et retours
paramétrables ; bonne qualité, reliés cartonnés disponibles.\par
\textbf{Limites :} interface et support en anglais ; il faut en général
fournir \textbf{son propre ISBN} ; moins adapté à un auteur débutant.
\end{fiche}

\begin{fiche}{Bookelis}{bookelis.com}
\textbf{Plateforme française d’autoédition.} Publication gratuite
possible, ISBN offert, impression à la demande et petits tirages,
diffusion en ligne et commande en librairie ; services payants à la
carte (correction, couverture, promotion).\par
\textbf{Points forts :} interlocuteur français, outils simples,
accompagnement.\par
\textbf{Limites :} les services additionnels peuvent vite alourdir la
facture ; comparez les marges par exemplaire.
\end{fiche}

\begin{fiche}{Les autres plateformes à connaître}{}
\begin{itemize}
  \item \textbf{TheBookEdition} (thebookedition.com) : autoédition et
        impression en France, boutique intégrée, petits tirages.
  \item \textbf{Librinova} (librinova.com) : autoédition accompagnée,
        avec un service d’agent qui repère les titres qui se vendent pour
        les proposer à des éditeurs.
  \item \textbf{Lulu} (lulu.com) : plateforme américaine historique,
        large choix de formats, réseau d’impression international.
  \item \textbf{Les éditeurs « à compte d’auteur »} : prudence. Si l’on
        vous demande de payer pour être « publié », lisez le contrat de
        très près et comparez avec ce que propose une plateforme d’IAD.
\end{itemize}
\end{fiche}

\subsection{Les imprimeurs numériques}

Pour une petite maison d’édition, une autre option consiste à travailler
directement avec un \textbf{imprimeur numérique} qui fabrique à l’unité
ou en très petites séries, et à gérer soi-même la vente (salons, site,
distributeur). Plusieurs imprimeurs français et européens proposent ce
service (par exemple Lightning Source à Maurepas ou le groupe CPI).
L’intérêt : un meilleur coût unitaire et plus de choix de papiers. La
contrainte : c’est à vous de gérer la logistique et la distribution.

\subsection{L’interprofession : les passages obligés}

\begin{center}
\small
\begin{tabularx}{\linewidth}{P{3.3cm}L}
\entete \textbf{Organisme} & \textbf{Rôle pour vous} \\
\textbf{AFNIL} & attribue gratuitement les ISBN aux éditeurs français \\
\textbf{BnF} & reçoit le dépôt légal ; ses notices alimentent les catalogues \\
\textbf{Dilicom} & portail par lequel les libraires commandent aux
  distributeurs ; y être référencé est indispensable pour la librairie \\
\textbf{Electre} & base bibliographique de référence des libraires et
  bibliothécaires \\
\textbf{CNL} & Centre national du livre : aides aux éditeurs et aux
  auteurs (conditions d’éligibilité à vérifier, l’autoédition en est
  souvent exclue) \\
\textbf{Agences régionales du livre} & aides, formations, réseau,
  salons dans votre région \\
\textbf{SGDL} & Société des gens de lettres : informations juridiques
  pour les auteurs, contrats \\
\textbf{Urssaf artistes-auteurs} & déclaration et cotisations des
  artistes-auteurs \\
\bottomrule
\end{tabularx}
\end{center}

\subsection{Comparatif rapide}

\begin{center}
\footnotesize
\begin{tabularx}{\linewidth}{P{2.1cm}P{2.4cm}LLP{2.4cm}}
\entete \textbf{Plateforme} & \textbf{Mise en ligne} & \textbf{ISBN} &
  \textbf{Librairies françaises} & \textbf{Idéal pour} \\
Amazon KDP & gratuite & gratuit (Amazon) ou le vôtre & quasi absent & ventes Amazon \\
\rowcolor{plCreme} BoD & offre gratuite ou forfaits & fourni ou le vôtre &
  oui, via Dilicom & auteurs \& éditeurs visant librairie + en ligne \\
IngramSpark & gratuite & le vôtre & selon réseau et conditions & éditeurs, international \\
Bookelis & gratuite + services & fourni & oui & auteurs accompagnés \\
TheBookEdition & gratuite + services & fourni ou le vôtre & limité & impression en France, petits tirages \\
Lulu & gratuite & fourni ou le vôtre & limité & formats variés, international \\
\bottomrule
\end{tabularx}
\end{center}
\source{Comparatif indicatif, établi en novembre 2026. Les offres
évoluent : vérifiez les conditions à jour sur chaque site.}
